\documentclass{amsart}
% For dealing with references we use the comment environment
\usepackage{verbatim}
\newenvironment{reference}{\comment}{\endcomment}
%\newenvironment{reference}{}{}
\newenvironment{slogan}{\comment}{\endcomment}
\newenvironment{history}{\comment}{\endcomment}

% For commutative diagrams we use Xy-pic
\usepackage[all]{xy}

% We use 2cell for 2-commutative diagrams.
\xyoption{2cell}
\UseAllTwocells

% We use multicol for the list of chapters between chapters
\usepackage{multicol}

% This is generally recommended for better output
\usepackage{lmodern}
\usepackage[T1]{fontenc}

% For cross-file-references

% Package for hypertext links:
\usepackage{hyperref}

% For any local file, say "hello.tex" you want to link to please

% Theorem environments.
%
\theoremstyle{plain}
\newtheorem{theorem}[subsection]{Theorem}
\newtheorem{proposition}[subsection]{Proposition}
\newtheorem{lemma}[subsection]{Lemma}

\theoremstyle{definition}
\newtheorem{definition}[subsection]{Definition}
\newtheorem{example}[subsection]{Example}
\newtheorem{exercise}[subsection]{Exercise}
\newtheorem{situation}[subsection]{Situation}

\theoremstyle{remark}
\newtheorem{remark}[subsection]{Remark}
\newtheorem{remarks}[subsection]{Remarks}

\numberwithin{equation}{subsection}

% Macros
%
\def\lim{\mathop{\mathrm{lim}}\nolimits}
\def\colim{\mathop{\mathrm{colim}}\nolimits}
\def\Spec{\mathop{\mathrm{Spec}}}
\def\Hom{\mathop{\mathrm{Hom}}\nolimits}
\def\Ext{\mathop{\mathrm{Ext}}\nolimits}
\def\SheafHom{\mathop{\mathcal{H}\!\mathit{om}}\nolimits}
\def\SheafExt{\mathop{\mathcal{E}\!\mathit{xt}}\nolimits}
\def\Sch{\mathit{Sch}}
\def\Mor{\mathop{\mathrm{Mor}}\nolimits}
\def\Ob{\mathop{\mathrm{Ob}}\nolimits}
\def\Sh{\mathop{\mathit{Sh}}\nolimits}
\def\NL{\mathop{N\!L}\nolimits}
\def\CH{\mathop{\mathrm{CH}}\nolimits}
\def\proetale{{pro\text{-}\acute{e}tale}}
\def\etale{{\acute{e}tale}}
\def\QCoh{\mathit{QCoh}}
\def\Ker{\mathop{\mathrm{Ker}}}
\def\Im{\mathop{\mathrm{Im}}}
\def\Coker{\mathop{\mathrm{Coker}}}
\def\Coim{\mathop{\mathrm{Coim}}}

% Boxtimes
%
\DeclareMathSymbol{\boxtimes}{\mathbin}{AMSa}{"02}

%
% Macros for moduli stacks/spaces
%
\def\QCohstack{\mathcal{QC}\!\mathit{oh}}
\def\Cohstack{\mathcal{C}\!\mathit{oh}}
\def\Spacesstack{\mathcal{S}\!\mathit{paces}}
\def\Quotfunctor{\mathrm{Quot}}
\def\Hilbfunctor{\mathrm{Hilb}}
\def\Curvesstack{\mathcal{C}\!\mathit{urves}}
\def\Polarizedstack{\mathcal{P}\!\mathit{olarized}}
\def\Complexesstack{\mathcal{C}\!\mathit{omplexes}}
% \Pic is the operator that assigns to X its picard group, usage \Pic(X)
% \Picardstack_{X/B} denotes the Picard stack of X over B
% \Picardfunctor_{X/B} denotes the Picard functor of X over B
\def\Pic{\mathop{\mathrm{Pic}}\nolimits}
\def\Picardstack{\mathcal{P}\!\mathit{ic}}
\def\Picardfunctor{\mathrm{Pic}}
\def\Deformationcategory{\mathcal{D}\!\mathit{ef}}

\setlength{\emergencystretch}{3em}
\begin{document}
\title[Stacks Project, Tag 092Z]{415 --- Olivier rigidity theorem for weakly etale local maps}
\maketitle
\noindent Copyright (C) 2005--2025 Johan de Jong. Stacks Project authors. Source: \href{https://stacks.math.columbia.edu/tag/092Z}{Tag 092Z}.

\noindent Editorial difficulty note (not author text). Difficulty H3: The proof shows every fibre has exactly one point with unchanged residue field. A hypothetical splitting yields an idempotent after an algebraic field extension, while integral closure over the strictly henselian base forces a local subring that excludes that idempotent; the map is then shown to be a faithfully flat epimorphism..

\noindent Editorial provenance note (not author text). History: excerpt from revision \texttt{a0444\allowbreak 6e57e\allowbreak c1fbc\allowbreak 252a8\allowbreak 71afc\allowbreak ec775\allowbreak 2fb28\allowbreak 07b14}, more-algebra.tex, lines 30787--30829. Reference presentation was adapted; original-author context and core support blocks were retained or added. The mathematical wording is unchanged. Licensed under GNU FDL 1.2 or later; no invariant sections or cover texts. See ../licenses/COPYING. Context blocks reproduce author definitions and supporting results; they are not additional counted problems.

\begin{theorem}[Olivier]
\label{theorem-olivier}
Let $A \to B$ be a local homomorphism of local rings.
If $A$ is strictly henselian and $A \to B$ is weakly \'etale, then
$A = B$.
\end{theorem}

\begin{proof}
We will show that for all $\mathfrak p \subset A$ there is a unique
prime $\mathfrak q \subset B$ lying over $\mathfrak p$ and
$\kappa(\mathfrak p) = \kappa(\mathfrak q)$.
This implies that $B \otimes_A B \to B$ is bijective on spectra
as well as surjective and flat. Hence it is an isomorphism
for example by the description of pure ideals in
Algebra, Lemma \href{https://stacks.math.columbia.edu/tag/04PU}{lemma pure open closed specializations (Tag 04PU)}.
Hence $A \to B$ is a faithfully flat epimorphism of rings. We get
$A = B$ by
Algebra, Lemma \href{https://stacks.math.columbia.edu/tag/04VU}{lemma faithfully flat epimorphism (Tag 04VU)}.

\medskip\noindent
Note that the fibre ring $B \otimes_A \kappa(\mathfrak p)$
is a colimit of \'etale extensions of $\kappa(\mathfrak p)$ by
Lemmas \href{https://stacks.math.columbia.edu/tag/092H}{lemma base change weakly etale (Tag 092H)} and
\href{https://stacks.math.columbia.edu/tag/092Q}{lemma absolutely flat over field (Tag 092Q)}.
Hence, if there exists more than one prime lying over $\mathfrak p$
or if $\kappa(\mathfrak p) \not = \kappa(\mathfrak q)$ for some $\mathfrak q$,
then $B \otimes_A L$ has a nontrivial idempotent for some (separable)
algebraic field extension $L/\kappa(\mathfrak p)$.

\medskip\noindent
Let $L/\kappa(\mathfrak p)$ be an algebraic field extension.
Let $A' \subset L$ be the integral closure of $A/\mathfrak p$ in $L$.
By Lemma \href{https://stacks.math.columbia.edu/tag/092X}{lemma integral over henselian (Tag 092X)}
we see that $A'$ is a strictly henselian local ring
whose residue field is a purely inseparable extension of the residue
field of $A$. Thus $B \otimes_A A'$ is a local ring by
Algebra, Lemma \href{https://stacks.math.columbia.edu/tag/092Y}{lemma local tensor with integral (Tag 092Y)}.
On the other hand, $B \otimes_A A'$ is integrally closed in
$B \otimes_A L$ by Lemma \href{https://stacks.math.columbia.edu/tag/092W}{lemma normality goes up (Tag 092W)}.
Since $B \otimes_A A'$ is local, it follows that the ring
$B \otimes_A L$ does not have nontrivial
idempotents which is what we wanted to prove.
\end{proof}
\end{document}
